\documentclass[11pt]{article}

\usepackage[T1]{fontenc}
\usepackage{amsmath,amssymb,amsthm}
\usepackage[margin=1in]{geometry}
\usepackage[hidelinks]{hyperref}

\newtheorem{theorem}{Theorem}
\newtheorem{remark}{Remark}
\newcommand{\AO}{\mathcal{AO}}
\newcommand{\psia}{\psi}

\title{A Modulo-Four Transfer Invariant for Acyclic Orientations of Width-Three Grid Strips}
\author{Xiaojian Zeng\\
PBC School of Finance, Tsinghua University, Beijing 100190, China\\
\texttt{xiaojianzeng@foxmail.com}}
\date{}

\begin{document}
\maketitle

\begin{abstract}
Let $\psia(G)$ denote the number of acyclic orientations of a finite simple
graph $G$.  Mühlherr and Poullot conjectured that
$\psia(P_n\times P_m)\equiv2\;(\bmod 4)$ whenever $m,n\geq3$ are odd,
after verifying three width-three instances computationally.  We prove the
complete width-three subfamily, and slightly more:
$\psia(P_3\times P_m)\equiv2\;(\bmod 4)$ for every $m\geq2$.
The proof is an elementary two-state transfer for the chromatic polynomial.
At $q=-1$, its state reaches $(2,0)$ modulo four after one extension and then
remains fixed.  The broader odd-by-odd conjecture is not addressed.
\end{abstract}

\section{Context and statement}

For the path $P_m$ on $m$ vertices, the Cartesian product
$P_3\times P_m$ is the usual $3$-by-$m$ rectangular grid.  If
$\chi_G(q)$ is the chromatic polynomial, Stanley's evaluation gives
\begin{equation}\label{eq:stanley}
  \psia(G)=(-1)^{|V(G)|}\chi_G(-1).
\end{equation}
This identity turns the orientation count into an exact negative evaluation
of a coloring invariant~\cite{Stanley1973}.

Mühlherr and Poullot reported
$\psia(P_3\times P_m)\equiv2\;(\bmod 4)$ for $m=3,5,7$ as computer
experiments and included all odd-by-odd grids in their Conjecture~1(1)
~\cite[Example~3.13(8) and Conjecture~1(1)]{MuehlherrPoullot2026}.
Transfer matrices for chromatic polynomials of square-lattice strips of fixed
width and arbitrary length have a substantial precedent; in particular,
Salas and Sokal developed the Fortuin--Kasteleyn transfer framework for free
strips of width at most eight~\cite{SalasSokal2001}.  Our proof needs only two
proper-coloring equality types at the exposed column and then reduces the
resulting matrix modulo four.

Coloring invariants also carry structural information beyond the numerical
evaluation used here.  For example, Zeng used the chromatic symmetric function
to reconstruct a multiplicity-isolated class of proper trees
~\cite{Zeng2026MultiplicityArxiv}.  That reconstruction problem is separate
from the strip enumeration below; both use coloring-derived graph invariants,
but for different questions.

\begin{theorem}\label{thm:main}
For every integer $m\geq2$,
\[
  \psia(P_3\times P_m)\equiv2\;(\bmod 4).
\]
Consequently, for every integer $k\geq1$,
\[
  \psia(P_3\times P_{2k+1})\equiv2\;(\bmod 4).
\]
\end{theorem}

The lower endpoint in the stronger assertion is necessary: $P_3$ has four
acyclic orientations, so the conclusion fails at $m=1$.

\section{The two-state transfer}

We first take an integer $q\geq3$ and count proper $q$-colorings.  A proper
coloring of the last three-vertex column has exactly one of the following two
equality types:
\[
  A=(a,b,a)\quad(a\neq b),
  \qquad
  B=(a,b,c)\quad(a,b,c\text{ pairwise distinct}).
\]
Let $A_m(q)$ and $B_m(q)$ count the colorings of the first $m$ columns whose
last column has type $A$ and $B$, respectively, and write the row vector
$v_m=(A_m,B_m)$.  For one column,
\begin{equation}\label{eq:initial}
  v_1=\bigl(q(q-1),q(q-1)(q-2)\bigr).
\end{equation}

Fix representatives $(0,1,0)$ and $(0,1,2)$ for the old types.  Direct
extension counting gives the old-type-by-new-type matrix
\begin{equation}\label{eq:matrix}
M(q)=
\left(\begin{array}{cc}
q^2-3q+3 &(q-2)(q^2-4q+5)\\
q^2-4q+5&q^3-6q^2+14q-13
\end{array}\right).
\end{equation}
Here every new vertex must differ both from its horizontal predecessor and
from its vertical neighbor or neighbors in the new column.  We include the
four counts for completeness.

For $A\to A$, separate the case in which the repeated new endpoint color is
$1$.  The count is
\[
 (q-1)+(q-2)^2=q^2-3q+3.
\]
For $B\to A$, the same separation gives
\[
 (q-1)+(q-3)(q-2)=q^2-4q+5.
\]
For $A\to B$, separate according as the new middle color is $0$ or is neither
$0$ nor $1$.  The number of extensions is
\[
 (q-1)(q-2)+(q-2)^2(q-3)
   =(q-2)(q^2-4q+5).
\]
Finally, for $B\to B$, start with all $q(q-1)(q-2)$ ordered distinct triples
and exclude equality with the corresponding entry of $(0,1,2)$.  Inclusion--
exclusion yields
\[
 q(q-1)(q-2)-3(q-1)(q-2)+3(q-2)-1
 =q^3-6q^2+14q-13.
\]
Color symmetry makes each extension number depend only on the old type.
Therefore
\begin{equation}\label{eq:transfer}
  v_{m+1}=v_mM(q),
  \qquad
  \chi_{P_3\times P_m}(q)=A_m(q)+B_m(q).
\end{equation}
Equations~\eqref{eq:initial}--\eqref{eq:transfer} hold for every integer
$q\geq3$.  Both sides are polynomials in $q$, so they are polynomial
identities and may legitimately be specialized at $q=-1$.

\section{The modulo-four fixed state}

At the required specialization,
\[
 M(-1)=\left(\begin{array}{cc}7&-30\\10&-34\end{array}\right)
 \equiv\left(\begin{array}{cc}3&2\\2&2\end{array}\right)\;(\bmod 4),
 \qquad
 v_1(-1)=(2,-6)\equiv(2,2)\;(\bmod 4).
\]
One multiplication gives
\[
  v_2(-1)\equiv(2,0)\;(\bmod 4),
\]
and this state is fixed because
\[
  (2,0)M(-1)\equiv(2,0)\;(\bmod 4).
\]
Induction in~\eqref{eq:transfer} now gives, for every $m\geq2$,
\[
 \chi_{P_3\times P_m}(-1)
 =A_m(-1)+B_m(-1)\equiv2\;(\bmod 4).
\]
Equation~\eqref{eq:stanley} changes this value by at most a sign, and
$-2\equiv2\;(\bmod 4)$.  This proves Theorem~\ref{thm:main}.

\begin{remark}
The theorem settles only the width-three subfamily of the odd-by-odd grid
conjecture.  No conclusion is claimed here for a fixed width at least five.
For the original odd-length grid, the number of edges is even, so the modular
criterion of Mühlherr and Poullot also gives the non-Hamiltonicity consequence
for the associated graph of acyclic orientations
~\cite[Lemma~3.2]{MuehlherrPoullot2026}.
\end{remark}

\section*{Exact checks and assistance disclosure}

The symbolic transfer above proves the quantified statement.  As checks, an
exact Python implementation enumerated all four transition types for
$q=3,\ldots,12$, compared the complete two-state chromatic polynomial with an
independent Fortuin--Kasteleyn connectivity transfer for $m=1,\ldots,9$, and
verified the fixed residue state.  At the first length beyond those reported
in~\cite{MuehlherrPoullot2026}, two independent exact programs---one using
deletion--contraction and one using the full chromatic polynomial---both gave
\[
 \psia(P_3\times P_9)=514302841174\equiv2\;(\bmod 4).
\]
These finite computations are corroborating checks, not premises of the
induction.

OpenAI Codex assisted with proof development, exact-check scripting, and
manuscript preparation.  The universal conclusion is supported by the
displayed symbolic argument; the accompanying scripts use exact integer and
polynomial arithmetic and introduce no floating-point assumptions.

\paragraph{Data and code availability.}
The accompanying supplement contains the width-three transfer certificate,
an exact deletion--contraction counter, an independent Fortuin--Kasteleyn
evaluator, and their saved outputs.

\bibliographystyle{plain}
\bibliography{references}

\end{document}
